\section{Related Work}

\subsection{Traditional Light Field Depth Estimation Methods}

Traditional light field depth estimation methods predominantly rely on physical cues and geometric constraints derived from Epipolar Plane Images (EPIs) . The foundational premise of these approaches is that a 3D point projects onto a straight line in the EPI, where the slope of this line is inversely proportional to the depth of the point. Building upon this linear structure, Wang et al. introduced an angular cone constraint theory to address depth discontinuities and occlusions. By modeling the viewing frustum of a light field camera as a continuous angular cone, their method formulates depth estimation as an optimization problem that enforces geometric consistency across adjacent EPI lines. This angular cone formulation effectively suppresses spurious depth hypotheses near object boundaries, yielding considerable improvements in preserving sharp depth edges compared to basic local slope estimators.

While spatial-domain slope calculation performs adequately in well-textured regions, it struggles in textureless or repetitive areas. To compensate for these limitations, early research integrated defocus cues and frequency-domain analyses. Tao et al. combined defocus responses with correspondence cues, leveraging the focal stack to resolve ambiguities where EPI slopes are ill-defined. In the frequency domain, methods employing 2D Discrete Fourier Transform (2D-DFT) or high-frequency filtering estimate EPI slopes by identifying the dominant orientation of the spectral energy . More recently, to capture complex local structures, Cui et al. proposed an X-pattern feature analysis that extracts cross-shaped structural signatures from EPIs. By analyzing the intersection and divergence of these X-patterns, their approach enhances depth discrimination in regions with intricate textures, demonstrating that explicit geometric pattern extraction can substantially improve robustness over simple gradient-based slope calculations.

Despite their theoretical elegance, these traditional methods exhibit substantial limitations when applied to unconstrained real-world environments. A primary deficiency is their strict reliance on the Lambertian reflectance assumption. The EPI angular cone constraints and linear structural models assume that pixel radiance remains invariant across viewpoints. Consequently, these geometric constraints fail completely on non-Lambertian surfaces, such as specular or scattering materials, where the physical laws governing light transport are inherently viewpoint-dependent, destroying the linear EPI structures . Under discrete and low angular resolutions (e.g., $9 \times 9$ sub-aperture images), traditional 2D-DFT frequency-domain analyses and high-frequency filtering techniques suffer from spectral aliasing and insufficient angular sampling. These frequency-based methods become highly susceptible to slope ambiguity in complex texture regions, lacking the robustness required for reliable depth recovery. These degradations indicate a pressing need to reveal the underlying physical and geometric causes and to extract more robust angular geometric features. This motivates our proposed three-layer angular signal decomposition and non-Lambertian detection mechanism, which explicitly models these physical deviations, alongside a unified dual-mask physical model designed to restore geometric consistency across heterogeneous material domains.

\subsection{Deep Learning-based General Light Field Depth Estimation Methods}

Following the limitations of handcrafted physical cues, deep learning-based approaches have emerged as the dominant paradigm for light field depth estimation. Shin et al. introduced EPINet , a foundational convolutional neural network (CNN) architecture that reorganizes 4D light field data into 2D sub-aperture image arrays or epipolar plane image (EPI)s (EPIs). By employing 2D convolutional kernels, EPINet simultaneously extracts spatial and angular features. To enhance directional sensitivity, subsequent variants such as EPINet4Dir process horizontal, vertical, and diagonal EPIs independently, fusing multi-directional features to sharpen depth boundaries. These CNN-based architectures demonstrate strong performance in well-textured Lambertian regions, substantially reducing matching ambiguities inherent in traditional methods. Nevertheless, their fixed receptive fields limit their effectiveness when handling large disparities and complex occlusions.

To overcome the constraints of fixed receptive fields and better capture global context, researchers have developed more sophisticated feature extraction and decoupling mechanisms. Wang et al. proposed DistgSSR , which adopts a disentangled strategy to separate light field features into spatial and angular components. By designing specific spatial-angular decoupled convolutions, this network independently optimizes spatial details and angular consistency, thereby improving depth regression accuracy while preserving high-frequency textures. Similarly, LF-DFN incorporates deformable convolutions and dynamic feature aggregation, enabling the network to adaptively adjust sampling locations based on local disparities. This geometry-aware feature extraction effectively mitigates depth artifacts in occluded regions, yielding considerable accuracy improvements on standard light field datasets.

With the advent of vision Transformers, recent studies have explored self-attention mechanisms to model long-range dependencies in light field data. Transformer-based architectures replace local convolutions with global self-attention, directly capturing spatial-angular correlations across the entire 4D light field. By computing attention weights between different viewpoints and spatial positions, these methods integrate global geometric cues more efficiently, exhibiting stronger robustness in weakly textured and repetitive pattern regions compared to pure CNN architectures. Although these advanced models continuously improve performance metrics in ideal scenarios, their high computational complexity remains a concern, and they still inherently rely on the underlying photo-consistency assumption.

Despite the considerable progress achieved by the aforementioned deep learning methods on standard benchmarks, they share a major limitation: most networks implicitly assume the scene consists of purely Lambertian surfaces, lacking explicit modeling of material reflectance properties. This assumption leads to severe feature confusion when processing heterogeneous materials in mixed scenes. In real-world environments, non-Lambertian regions violate photo-consistency across viewpoints, disrupting the learned feature representations. Specifically, a single network architecture struggles to simultaneously accommodate the complex texture recovery required for Lambertian regions and the depth regression needed for non-Lambertian specular areas. Existing methods lack an adaptive feature routing mechanism to dynamically adjust processing strategies based on local material attributes. Consequently, introducing a mask-guided dual-branch or routing architecture is necessary to isolate and process different material regions independently. This deficiency directly motivates our approach. By designing a three-layer angular signal decomposition and non-Lambertian detection mechanism based on geometric features, alongside a unified dual-mask physical model and multi-domain balanced training strategy, our method explicitly decouples Lambertian and non-Lambertian signals, enabling robust depth estimation across complex material interactions.

\subsection{Deep Learning-based Non-Lambertian and Mask-Guided Depth Estimation Methods}

To mitigate the matching ambiguities caused by non-Lambertian reflectance and complex occlusions, recent advancements have incorporated mask-guided mechanisms and multi-branch architectures into light field depth estimation. Instead of treating all angular pixels equally, these methods attempt to identify and isolate non-Lambertian regions, assigning them lower confidence weights during depth regression.

Early mask-guided approaches primarily relied on heuristic rules or independent binary classifiers to detect non-Lambertian surfaces. For instance, Zhang et al. utilized the variance of epipolar plane image (EPI) (EPI) gradients and photo-consistency errors to generate a specular mask, which subsequently re-weighted the cost volume during depth inference. While this strategy reduces the influence of specular highlights, the heuristic thresholds are highly sensitive to noise and texture variations. To improve detection accuracy, subsequent works employed black-box neural networks for mask generation. Chen et al. trained a separate segmentation network to predict non-Lambertian masks from central sub-aperture images, feeding the masks into the main depth estimation network as spatial attention maps. Although this data-driven approach yields higher mask precision on synthetic datasets, it lacks physical interpretability and struggles to generalize to unseen material properties, as the mask generation is entirely decoupled from the underlying light transport physics.

To better accommodate the heterogeneous nature of mixed-material scenes, researchers have developed dual-branch architectures that process Lambertian and non-Lambertian features separately. Huang et al. proposed a dual-branch network where one branch extracts global geometric context from EPIs, while the other focuses on local texture and specular features from sub-aperture images. The features from both branches are fused using a cross-attention mechanism guided by a predicted occlusion and non-Lambertian mask. This separation allows the network to preserve sharp depth boundaries near specular regions. However, processing distinct features in isolated branches often fragments the global context, and the fixed fusion weights lack adaptive feature routing within a unified model, limiting the network's capacity to handle gradual material transitions.

Recognizing the complementary nature of different optical cues, multi-task models have introduced defocus-based auxiliary heads to assist in non-Lambertian identification. Liu et al. [1] integrated a defocus estimation branch alongside the primary correspondence-based depth branch. Since non-Lambertian surfaces exhibit distinct inconsistencies between defocus cues (which depend on the lens aperture and focus distance) and correspondence cues (which rely on viewpoint shifts), the discrepancy between the two auxiliary outputs serves as a strong indicator for specular and translucent regions. By jointly optimizing depth and defocus maps, this method effectively suppresses depth artifacts on reflective surfaces. Despite these improvements, the auxiliary heads significantly increase computational overhead, and the multi-task learning process is highly susceptible to the severe data scarcity of real-world non-Lambertian samples, leading to sub-optimal convergence in complex environments.

Despite the progress made by mask-guided and multi-branch methods, several significant limitations persist. First, existing non-Lambertian detection mechanisms predominantly depend on simple heuristic rules or black-box neural networks. These approaches lack a physics-based angular signal decomposition—such as decoupling ambient illumination, parallax, and material residuals—which substantially restricts the precision and physical consistency of the generated masks. Second, dual-branch architectures often fragment the global context and lack adaptive feature routing within a unified framework. More importantly, current methods lack systematic training strategies, such as domain-balanced sampling, to address the severe scarcity of non-Lambertian data. Consequently, these models exhibit insufficient generalization capabilities when deployed in real-world complex urban and mixed-material scenes, where diverse reflectance properties coexist.

Although existing mask-guided and multi-branch methods have alleviated matching ambiguities to some extent, they typically treat non-Lambertian detection and depth estimation as decoupled processes and lack explicit physical modeling of angular signals, leading to suboptimal performance in complex reflective and occluded scenes. Inspired by these observations, we propose a unified dual-mask physical modeling approach that addresses the aforementioned limitations by incorporating a geometry-aware three-layer angular signal decomposition for accurate non-Lambertian detection and a multi-domain balanced training strategy to robustly optimize depth regression. Building upon this conceptual overview, we now detail the complete pipeline, beginning with the overall architecture and core mathematical formulations of our unified dual-mask physical model.